\documentclass[10pt,a4paper]{amsart}
\usepackage[margin=19mm]{geometry}
\usepackage{amsmath,amssymb,mathtools}
\usepackage[scr=rsfs,cal=euler]{mathalfa}
\usepackage{xfrac}
\usepackage[unicode,hidelinks,bookmarks=false]{hyperref}
\newcommand{\ie}{i.e.\ }
\newcommand{\cf}{cf.\ }
\newcommand{\resp}{respectively\ }
\newcommand{\Subsection}{\subsection}
\providecommand{\nobreakdash}{\nobreak\hbox{-}}
\setlength{\emergencystretch}{2em}
\multlinegap0pt
\usepackage{amssymb}
\usepackage{mathtools}
\usepackage[shortlabels]{enumitem}
\newtheorem{proposition}{Proposition}
\title{Inhomogeneous Approximation by Sums of Roots}
\author{Samuel Korsky}
\date{Source: arXiv:2605.27233v1 (CC BY 4.0)}
\begin{document}
\maketitle

\noindent\emph{Source and licence.} Inhomogeneous Approximation by Sums of Roots. Version arXiv:2605.27233v1 (CC BY 4.0), distributed by arXiv under the Creative Commons Attribution 4.0 International licence, http://creativecommons.org/licenses/by/4.0/, as recorded in the arXiv licence metadata for this exact version and shown on the abstract page https://arxiv.org/abs/2605.27233v1. Full licence notice: \texttt{licenses/arxiv-2605.27233-CC-BY-4.0.txt}.\par\smallskip

\noindent\textit{Editorial scope.} Counted unit: the article's proposition prop:examples (source0.tex lines 477--488). The article's complete original proof of it is delivered with it (source0.tex lines 490--559, one \texttt{proof} environment of 70 lines). No mathematics is written, repaired or re-derived: the statement and the proof are the article's own lines, in the article's own notation, and no retained line is substituted or re-flowed.  Retained uncounted as the source-local context the proof needs: lines 450--471.  Nothing was omitted from any retained span, so the package carries no omission record.  No retained line contains a citation, so the wrapper carries no bibliography and no bibliography entry.  No retained line contains a figure environment, an image-inclusion command or drawing code, so no image is delivered and the package declares no asset dependency.  Editorial formatting only, in addition: an AMS \texttt{amsart} preamble that restates the article's own declarations and macros used by the retained lines, namely the environments \texttt{proposition} on separate counters, together with the standard packages those lines need.  The retained environments print in this document as Proposition 1 in order of appearance, whereas the article prints the same items with the numbering of its own full document.  The complete native source of the article as distributed in the arXiv e-print for this exact version is bundled unchanged as \texttt{sources/201476/source0.tex}.
\bigskip
\noindent
Let
\[
  C_t(M):=\sqrt{M^2+t}
\]
and
\[
  P(M):=\sqrt{(M+1)^2+1}+\sqrt{(M-1)^2+1}.
\]
As $M\to\infty$,
\[
  C_t(M)=M+\frac{t}{2M}-\frac{t^2}{8M^3}
  +\frac{t^3}{16M^5}-\frac{5t^4}{128M^7}+O_t(M^{-9}),
  \tag{5.1}\label{eq:C-expansion}
\]
and
\[
  P(M)=2M+\frac1M+\frac{3}{4M^3}-\frac{3}{8M^5}
  -\frac{61}{64M^7}+O(M^{-9}).
  \tag{5.2}\label{eq:P-expansion}
\]

\begin{proposition}[Explicit square-root integer-target examples]\label{prop:examples}
For $k=2,3,4$, there are infinitely many $k$-tuples of positive integers $b_1,\ldots,b_k\ll M^2$ such that
\[
  0<\left\|\sum_{j=1}^k\sqrt{b_j}\right\|
  \ll_k M^{-(2k-1)}.
\]
Consequently, for every sufficiently large $N$,
\[
  g_{k,2}(N)\ll_k N^{-(k-1/2)}
  \qquad(k=2,3,4).
\]
\end{proposition}

\begin{proof}
For $k=2$, take
\[
  S_2(M)=\sqrt{M^2-1}+\sqrt{M^2+1}.
\]
By \eqref{eq:C-expansion},
\[
  S_2(M)=2M-\frac{1}{4M^3}+O(M^{-7}).
\]
Thus $0<\|S_2(M)\|\ll M^{-3}$ for all sufficiently large integers $M$.

\bigskip
\noindent
For $k=3$, take
\[
\begin{aligned}
  S_3(M)
  &:=3P(M)+2C_{-3}(M)\\
  &=\sqrt{9((M+1)^2+1)}
    +\sqrt{9((M-1)^2+1)}
    +\sqrt{4(M^2-3)}.
\end{aligned}
\]
Using \eqref{eq:C-expansion} and \eqref{eq:P-expansion},
\[
  S_3(M)=8M-\frac{9}{2M^5}+O(M^{-7}).
\]
Hence $0<\|S_3(M)\|\ll M^{-5}$.

\bigskip
\noindent
For $k=4$, take
\[
\begin{aligned}
  S_4(M)
  &:=99P(M)+108C_{-2}(M)+2C_9(M)\\
  &=\sqrt{99^2((M+1)^2+1)}
    +\sqrt{99^2((M-1)^2+1)}\\
  &\hspace{0.7cm}
    +\sqrt{108^2(M^2-2)}
    +\sqrt{2^2(M^2+9)}.
\end{aligned}
\]
Again using \eqref{eq:C-expansion} and \eqref{eq:P-expansion}, the coefficients of $M^{-1}$, $M^{-3}$, and $M^{-5}$ cancel:
\[
\begin{aligned}
  M^{-1}:&\quad 99-108+9=0,\\
  M^{-3}:&\quad 99\cdot\frac34+108\cdot\left(-\frac12\right)
  +2\cdot\left(-\frac{81}{8}\right)=0,\\
  M^{-5}:&\quad 99\cdot\left(-\frac38\right)+108\cdot\left(-\frac12\right)
  +2\cdot\frac{729}{16}=0.
\end{aligned}
\]
The first nonzero term is
\[
  99\left(-\frac{61}{64}\right)
  +108\left(-\frac58\right)
  +2\left(-\frac{32805}{128}\right)
  =-\frac{10791}{16}.
\]
Therefore
\[
  S_4(M)=308M-\frac{10791}{16M^7}+O(M^{-9}),
\]
and $0<\|S_4(M)\|\ll M^{-7}$.

\bigskip
\noindent
In all three cases the radicands are positive integers for all sufficiently large integer $M$, and their maximum is $\ll M^2$. Given any sufficiently large radicand bound $N$, choose $M=\lfloor c_k N^{1/2}\rfloor$ with $c_k>0$ small enough that all displayed radicands are at most $N$. This gives the claimed exponent $N^{-(k-1/2)}$ for every sufficiently large $N$.
\end{proof}

\end{document}
